\begin{proposition}[Smooth fibres of an algebraic reduction]\hypertarget{tag-01F4}{}\label{cdp:st:C.3.1}
\StatementRecord{cdp:st:C.3.1}{proved}{01F4}{\SeriesRef{cdp:st:C.3.7}, \SeriesRef{cdp:st:C.3.8}, \SeriesRef{cdp:st:C.3.9}, \SeriesRef{cdp:st:C.3.10}, \SeriesRef{cdp:st:C.3.6}, \SeriesRef{cdp:st:D.2.7}, \SeriesRef{cdp:st:C.3.12}; \cite[Theorem 3.1]{Peters1990}; \cite[p.~2]{Teleman2008}; \cite[Theorem 1]{BBK2019}}

Let $f:Y\to C$ be a holomorphic algebraic reduction of a connected
smooth compact complex threefold with $b_1(Y)=b_2(Y)=0$ and $a(Y)=1$.
The base is $C\simeq\PP^1$. Every smooth fibre $F$ has
\[
 c_1(K_F)=0\text{ in }H^2(F;\Q),\qquad
 c_2(F)=0\text{ in }H^4(F;\Z),\qquad
 \chi(F,\OO_F)=0,
\]
and contains no exceptional curve of the first kind. A very general
smooth fibre has $\kappa(F)\leq0$. The very general fibre's type belongs to the following
list: a complex two-torus, a bielliptic surface, a primary or secondary
Kodaira surface, a primary or secondary Hopf surface, or an Inoue
surface. Every listed type has Euler characteristic zero. This is a
preliminary necessary list. Statement \ref{cdp:st:C.3.7} below excludes the
bielliptic entry, and \ref{cdp:st:C.3.8} makes the very general Hopf fibre have
algebraic dimension zero. \ref{cdp:st:C.3.9} and \ref{cdp:st:C.3.10} exclude both Kodaira
entries. Statement \ref{cdp:st:C.3.6} realizes both primary and secondary Hopf
entries, and \ref{cdp:st:D.2.7} realizes a fixed abelian surface as general
fibre. Statement \ref{cdp:st:C.3.12} excludes the Inoue entry. Thus the
final very general list consists exactly of two-tori and
primary or secondary Hopf surfaces of algebraic dimension zero.
\end{proposition}


\begin{proof}

\begin{enumerate}[leftmargin=*]
\item Connected fibres give $f_*\Q_Y=\Q_C$ and the degree-one edge injection
$H^1(C;\Q)\hookrightarrow H^1(Y;\Q)$. Therefore the genus of $C$ is
zero. The holomorphic tangent sequence on a smooth fibre, with the
last map induced by the differential of $f$, is
\[
0\longrightarrow T_F\longrightarrow T_Y|_F
 \xrightarrow{df} \OO_F\otimes T_{C,f(F)}\longrightarrow0.
\]
The normal factor is a trivial holomorphic line bundle. Thus
$K_F=K_Y|_F$ and $c_2(F)=c_2(Y)|_F$. Rational $H^2(Y)$ and $H^4(Y)$
vanish, and $H^4(F;\Z)=\Z$ has no torsion. These facts give the two
Chern-class assertions. Noether's formula gives $\chi(\OO_F)=0$.
An exceptional curve $E\simeq\PP^1$ with $E^2=-1$ would have
$K_F\cdot E=-1$ by adjunction, contradicting the rational vanishing of
$c_1(K_F)$.

\item Fix $m>0$ and put $G_m=f_*(K_Y^m)$. This is a coherent torsion-free
sheaf on the smooth curve $C$: a section killed by a nonzero base
function vanishes on a dense open subset of $Y$ and hence everywhere.
It is consequently a vector bundle. On the nonempty open set of
cohomological base change its fibre is $H^0(F,K_F^m)$. A rank at least
two would give a meromorphic map
\[
 Y\dashrightarrow \PP_C(G_m^*)
\]
whose restriction to a general fibre is defined by at least two
independent sections of a line bundle and has positive-dimensional
image. Its image has dimension at least two. The target is projective:
a vector bundle on a compact Riemann surface is algebraic, and its
projectivization is a projective variety. The image is therefore
projective and its rational functions pull back to a subfield of
$\mathcal M(Y)$ of transcendence degree at least two. This contradicts
$a(Y)=1$. Thus $\rk G_m\leq1$ for every $m$. Removing the countable
union of the exceptional base-change sets gives the claimed very
general-fibre assertion.

\item The Enriques--Kodaira classification is given in
\cite[Theorem 3.1]{Peters1990}. The surface is minimal by the
adjunction argument. Rational and ruled surfaces have a rational
curve on which the canonical class has nonzero degree, so they are
excluded. K3 and Enriques surfaces have Euler numbers $24$ and $12$,
respectively. Properly elliptic and general-type surfaces have
positive Kodaira dimension. The remaining K\"ahler possibilities are
tori and bielliptic surfaces; the remaining Kodaira-dimension-zero
non-K\"ahler possibilities are the two types of Kodaira surfaces.
For class VII, $b_1(F)=1$, and $e(F)=0$ gives $b_2(F)=0$. The
classification of this case gives Hopf and Inoue surfaces; its
precise scope, including finite covers, is stated in
\cite[p.~2]{Teleman2008} and \cite[Theorem 1]{BBK2019}.

\end{enumerate}
\end{proof}